\begin{center}
\textbf{Refroidir une boisson}
\end{center}

Une élève de terminale STI2D a acheté des glaçons ayant des aspects très variés : matière plastique, acier inoxidable, pierre de granit.

Elle s'interroge sur leur efficacité en comparaison des glaçons traditionnels constitués d'eau sous forme de glace.

Elle réalise l'expérience suivante : après avoir versé 200 g d'eau dans quatre verres munis chacun d'un thermomètre, elle ajoute dans chacun un glaçon de volume identique, tout juste sorti du congélateur.

Elle suit ensuite l'évolution de la température de l'eau de chaque verre en fonction du temps.

\medskip

On s'intéresse à $T_1$\,, la température de la boisson (en \textdegree{}C) refroidie avec les glaçons d'aspect plastique.

La température $T_1$ est modélisée par la fonction définie pour tout $t \in [0~;~+\infty[$ par :
\[T_1(t) = 11{,}5\e^{-0{,}006t} + 11{,}5.\]

On admet que $T_1$ est dérivable sur $[0~;~+\infty[$ et on note $T_1'$ sa dérivée.

On considère l'équation différentielle suivante :
\[(E)~:~y' + 0{,}006y = 0{,}069.\]

\begin{enumerate}[start=4]
\item Montrer que $T_1$ est la solution de $E$ vérifiant la condition $T_1(0) = 23$.

\item Déterminer la valeur de $\displaystyle \lim_{t\to +\infty} T_1(t)$.

\item En supposant que la température extérieure reste constante et égale à 23,0\,\textdegree{}C, estimer la valeur de la température de la boisson au bout de 3~h.

Indiquer si le modèle mathématique précédent peut être appliqué dans ce cas.
\end{enumerate}
